\documentclass[12pt]{article}
\usepackage[utf8]{inputenc}
\usepackage[T1]{fontenc}
\usepackage{hyperref}
\title{Upper Necaxa Totonac Language Resources for K-12 Curriculum Development}
\author{Compiled by Claude Sonnet 5 (claude-sonnet-5) for Andr\'as Kornai}
\date{}
\begin{document}
\maketitle

\section{Introduction}
Upper Necaxa Totonac (tku) is spoken by about 3,400 people in Puebla State, Mexico, in two main dialects. David Beck's documentation project produced a dictionary and some community materials.

\section{Primary Educational Resources}
The project produced two practical vocabularies in Spanish, distributed free in Upper Necaxa communities (described on the project site). The search results did not give a direct link to them. \url{https://www.artsrn.ualberta.ca/totonaco2/?p=112}

\section{Linguistic and Reference Materials}
\begin{itemize}
\item David Beck, \emph{Upper Necaxa Totonac Dictionary}, De Gruyter Mouton, 2011. It has about 9,000 entries and a short grammatical sketch. A DVD adds audio and interlinear examples. \url{https://www.degruyterbrill.com/document/doi/10.1515/9783110238235/html}
\item Catalogue record at the Austrian Academy of Sciences. \url{https://oeaw.ac.at/resources/Record/9783110238235}
\item Review in \emph{Estudios de Ling\"u\'istica Aplicada}. \url{https://ela.enallt.unam.mx/index.php/ela/article/download/443/430}
\item Beck's handouts. \url{https://sites.ualberta.ca/~dbeck/handouts.pdf}
\end{itemize}

\section{Religious and Community Resources}
None were found.

\section{Digital and Audio Resources}
The dictionary DVD and electronic version include audio. The project website says a fuller grammar and a collection of analysed texts are in progress, but gives no date. \url{https://www.artsrn.ualberta.ca/totonaco2/?p=112}

\section{Curriculum Development Context}
The language is not known to be taught in schools. The dictionary and the Spanish vocabularies are the best starting points for a curriculum.
\end{document}